\documentclass{article}
% Competition/workshop paper, non-anonymous (single-blind) camera-ready: the official style
% requires a track option together with [final]; [final] alone leaves the track name undefined.
\usepackage[sglblindworkshop,final]{neurips_2026}
\workshoptitle{Smart Buildings Challenge (SB26)}
\usepackage[utf8]{inputenc}
\usepackage[T1]{fontenc}
\usepackage{url}
\usepackage{booktabs}
\usepackage{amsfonts}
\usepackage{amsmath}
\usepackage{nicefrac}
\usepackage{microtype}
\usepackage{xcolor}
\usepackage{graphicx}
\usepackage{algorithm}
\usepackage{algpseudocode}
\usepackage{multirow}
\usepackage{hyperref}  % after float/algorithm so caption anchors are not duplicated
% AUTO-GENERATED by experiments/src/make_macros.py -- do not edit by hand
\newcommand{\NScen}{30}
\newcommand{\MetricOurs}{0.0083}
\newcommand{\MetricSdOurs}{0.0093}
\newcommand{\MetricOursFairbanks}{0.0175}
\newcommand{\MetricOursDenver}{0.0050}
\newcommand{\MetricOursAtlanta}{0.0055}
\newcommand{\MetricOursPhoenix}{0.0075}
\newcommand{\MetricOursSingapore}{0.0060}
\newcommand{\MetricOursBone}{0.0159}
\newcommand{\MetricOursBtwo}{0.0071}
\newcommand{\MetricOursBfour}{0.0052}
\newcommand{\MetricOursBfive}{0.0041}
\newcommand{\MetricOursBseven}{0.0071}
\newcommand{\MetricOursBeight}{0.0104}
\newcommand{\MetricOursBgena}{\textbf{??}}
\newcommand{\MetricOursBgenb}{\textbf{??}}
\newcommand{\DecompOursComfort}{0.0004}
\newcommand{\DecompOursComfortFairbanks}{0.0008}
\newcommand{\DecompOursComfortDenver}{0.0004}
\newcommand{\DecompOursComfortAtlanta}{0.0002}
\newcommand{\DecompOursComfortPhoenix}{0.0003}
\newcommand{\DecompOursComfortSingapore}{0.0002}
\newcommand{\DecompOursEnergy}{0.0040}
\newcommand{\DecompOursEnergyFairbanks}{0.0049}
\newcommand{\DecompOursEnergyDenver}{0.0021}
\newcommand{\DecompOursEnergyAtlanta}{0.0029}
\newcommand{\DecompOursEnergyPhoenix}{0.0056}
\newcommand{\DecompOursEnergySingapore}{0.0047}
\newcommand{\DecompOursCarbon}{0.0039}
\newcommand{\DecompOursCarbonFairbanks}{0.0118}
\newcommand{\DecompOursCarbonDenver}{0.0025}
\newcommand{\DecompOursCarbonAtlanta}{0.0024}
\newcommand{\DecompOursCarbonPhoenix}{0.0017}
\newcommand{\DecompOursCarbonSingapore}{0.0011}
\newcommand{\SecPerStepOurs}{0.058}
\newcommand{\PolicyMsOurs}{0.43}
\newcommand{\MetricOursVFour}{0.0092}
\newcommand{\MetricSdOursVFour}{0.0097}
\newcommand{\MetricOursVFourFairbanks}{0.0178}
\newcommand{\MetricOursVFourDenver}{0.0050}
\newcommand{\MetricOursVFourAtlanta}{0.0058}
\newcommand{\MetricOursVFourPhoenix}{0.0096}
\newcommand{\MetricOursVFourSingapore}{0.0078}
\newcommand{\MetricOursVFourBone}{0.0172}
\newcommand{\MetricOursVFourBtwo}{0.0082}
\newcommand{\MetricOursVFourBfour}{0.0056}
\newcommand{\MetricOursVFourBfive}{0.0044}
\newcommand{\MetricOursVFourBseven}{0.0079}
\newcommand{\MetricOursVFourBeight}{0.0120}
\newcommand{\MetricOursVFourBgena}{\textbf{??}}
\newcommand{\MetricOursVFourBgenb}{\textbf{??}}
\newcommand{\DecompOursVFourComfort}{0.0002}
\newcommand{\DecompOursVFourComfortFairbanks}{0.0005}
\newcommand{\DecompOursVFourComfortDenver}{0.0003}
\newcommand{\DecompOursVFourComfortAtlanta}{0.0001}
\newcommand{\DecompOursVFourComfortPhoenix}{0.0002}
\newcommand{\DecompOursVFourComfortSingapore}{0.0002}
\newcommand{\DecompOursVFourEnergy}{0.0048}
\newcommand{\DecompOursVFourEnergyFairbanks}{0.0050}
\newcommand{\DecompOursVFourEnergyDenver}{0.0023}
\newcommand{\DecompOursVFourEnergyAtlanta}{0.0032}
\newcommand{\DecompOursVFourEnergyPhoenix}{0.0073}
\newcommand{\DecompOursVFourEnergySingapore}{0.0062}
\newcommand{\DecompOursVFourCarbon}{0.0042}
\newcommand{\DecompOursVFourCarbonFairbanks}{0.0124}
\newcommand{\DecompOursVFourCarbonDenver}{0.0025}
\newcommand{\DecompOursVFourCarbonAtlanta}{0.0024}
\newcommand{\DecompOursVFourCarbonPhoenix}{0.0021}
\newcommand{\DecompOursVFourCarbonSingapore}{0.0015}
\newcommand{\SecPerStepOursVFour}{0.058}
\newcommand{\PolicyMsOursVFour}{0.46}
\newcommand{\MetricPidWater}{0.0550}
\newcommand{\MetricSdPidWater}{0.0454}
\newcommand{\MetricPidWaterFairbanks}{0.0791}
\newcommand{\MetricPidWaterDenver}{0.0379}
\newcommand{\MetricPidWaterAtlanta}{0.0306}
\newcommand{\MetricPidWaterPhoenix}{0.0801}
\newcommand{\MetricPidWaterSingapore}{0.0471}
\newcommand{\MetricPidWaterBone}{0.0871}
\newcommand{\MetricPidWaterBtwo}{0.0481}
\newcommand{\MetricPidWaterBfour}{0.0362}
\newcommand{\MetricPidWaterBfive}{0.0470}
\newcommand{\MetricPidWaterBseven}{0.0429}
\newcommand{\MetricPidWaterBeight}{0.0686}
\newcommand{\MetricPidWaterBgena}{\textbf{??}}
\newcommand{\MetricPidWaterBgenb}{\textbf{??}}
\newcommand{\DecompPidWaterComfort}{0.0349}
\newcommand{\DecompPidWaterComfortFairbanks}{0.0415}
\newcommand{\DecompPidWaterComfortDenver}{0.0322}
\newcommand{\DecompPidWaterComfortAtlanta}{0.0228}
\newcommand{\DecompPidWaterComfortPhoenix}{0.0498}
\newcommand{\DecompPidWaterComfortSingapore}{0.0282}
\newcommand{\DecompPidWaterEnergy}{0.0115}
\newcommand{\DecompPidWaterEnergyFairbanks}{0.0143}
\newcommand{\DecompPidWaterEnergyDenver}{0.0034}
\newcommand{\DecompPidWaterEnergyAtlanta}{0.0048}
\newcommand{\DecompPidWaterEnergyPhoenix}{0.0216}
\newcommand{\DecompPidWaterEnergySingapore}{0.0135}
\newcommand{\DecompPidWaterCarbon}{0.0086}
\newcommand{\DecompPidWaterCarbonFairbanks}{0.0233}
\newcommand{\DecompPidWaterCarbonDenver}{0.0024}
\newcommand{\DecompPidWaterCarbonAtlanta}{0.0031}
\newcommand{\DecompPidWaterCarbonPhoenix}{0.0087}
\newcommand{\DecompPidWaterCarbonSingapore}{0.0054}
\newcommand{\SecPerStepPidWater}{0.060}
\newcommand{\PolicyMsPidWater}{0.34}
\newcommand{\MetricPid}{0.0678}
\newcommand{\MetricSdPid}{0.0474}
\newcommand{\MetricPidFairbanks}{0.0759}
\newcommand{\MetricPidDenver}{0.0506}
\newcommand{\MetricPidAtlanta}{0.0432}
\newcommand{\MetricPidPhoenix}{0.0955}
\newcommand{\MetricPidSingapore}{0.0739}
\newcommand{\MetricPidBone}{0.0784}
\newcommand{\MetricPidBtwo}{0.0662}
\newcommand{\MetricPidBfour}{0.0394}
\newcommand{\MetricPidBfive}{0.0721}
\newcommand{\MetricPidBseven}{0.0594}
\newcommand{\MetricPidBeight}{0.0915}
\newcommand{\MetricPidBgena}{\textbf{??}}
\newcommand{\MetricPidBgenb}{\textbf{??}}
\newcommand{\DecompPidComfort}{0.0421}
\newcommand{\DecompPidComfortFairbanks}{0.0422}
\newcommand{\DecompPidComfortDenver}{0.0384}
\newcommand{\DecompPidComfortAtlanta}{0.0288}
\newcommand{\DecompPidComfortPhoenix}{0.0579}
\newcommand{\DecompPidComfortSingapore}{0.0433}
\newcommand{\DecompPidEnergy}{0.0143}
\newcommand{\DecompPidEnergyFairbanks}{0.0124}
\newcommand{\DecompPidEnergyDenver}{0.0064}
\newcommand{\DecompPidEnergyAtlanta}{0.0079}
\newcommand{\DecompPidEnergyPhoenix}{0.0253}
\newcommand{\DecompPidEnergySingapore}{0.0198}
\newcommand{\DecompPidCarbon}{0.0114}
\newcommand{\DecompPidCarbonFairbanks}{0.0214}
\newcommand{\DecompPidCarbonDenver}{0.0059}
\newcommand{\DecompPidCarbonAtlanta}{0.0065}
\newcommand{\DecompPidCarbonPhoenix}{0.0122}
\newcommand{\DecompPidCarbonSingapore}{0.0108}
\newcommand{\SecPerStepPid}{0.056}
\newcommand{\PolicyMsPid}{0.22}
\newcommand{\MetricFixedSp}{0.0655}
\newcommand{\MetricSdFixedSp}{0.0527}
\newcommand{\MetricFixedSpFairbanks}{0.0759}
\newcommand{\MetricFixedSpDenver}{0.0772}
\newcommand{\MetricFixedSpAtlanta}{0.0330}
\newcommand{\MetricFixedSpPhoenix}{0.1161}
\newcommand{\MetricFixedSpSingapore}{0.0253}
\newcommand{\MetricFixedSpBone}{0.0729}
\newcommand{\MetricFixedSpBtwo}{0.0758}
\newcommand{\MetricFixedSpBfour}{0.0355}
\newcommand{\MetricFixedSpBfive}{0.0432}
\newcommand{\MetricFixedSpBseven}{0.0921}
\newcommand{\MetricFixedSpBeight}{0.0735}
\newcommand{\MetricFixedSpBgena}{\textbf{??}}
\newcommand{\MetricFixedSpBgenb}{\textbf{??}}
\newcommand{\DecompFixedSpComfort}{0.0337}
\newcommand{\DecompFixedSpComfortFairbanks}{0.0371}
\newcommand{\DecompFixedSpComfortDenver}{0.0433}
\newcommand{\DecompFixedSpComfortAtlanta}{0.0231}
\newcommand{\DecompFixedSpComfortPhoenix}{0.0495}
\newcommand{\DecompFixedSpComfortSingapore}{0.0153}
\newcommand{\DecompFixedSpEnergy}{0.0176}
\newcommand{\DecompFixedSpEnergyFairbanks}{0.0140}
\newcommand{\DecompFixedSpEnergyDenver}{0.0179}
\newcommand{\DecompFixedSpEnergyAtlanta}{0.0052}
\newcommand{\DecompFixedSpEnergyPhoenix}{0.0428}
\newcommand{\DecompFixedSpEnergySingapore}{0.0081}
\newcommand{\DecompFixedSpCarbon}{0.0143}
\newcommand{\DecompFixedSpCarbonFairbanks}{0.0248}
\newcommand{\DecompFixedSpCarbonDenver}{0.0159}
\newcommand{\DecompFixedSpCarbonAtlanta}{0.0047}
\newcommand{\DecompFixedSpCarbonPhoenix}{0.0239}
\newcommand{\DecompFixedSpCarbonSingapore}{0.0019}
\newcommand{\SecPerStepFixedSp}{0.058}
\newcommand{\PolicyMsFixedSp}{0.24}
\newcommand{\MetricBangBang}{0.2245}
\newcommand{\MetricSdBangBang}{0.0593}
\newcommand{\MetricBangBangFairbanks}{0.1996}
\newcommand{\MetricBangBangDenver}{0.1830}
\newcommand{\MetricBangBangAtlanta}{0.2315}
\newcommand{\MetricBangBangPhoenix}{0.2367}
\newcommand{\MetricBangBangSingapore}{0.2718}
\newcommand{\MetricBangBangBone}{0.2116}
\newcommand{\MetricBangBangBtwo}{0.2440}
\newcommand{\MetricBangBangBfour}{0.1861}
\newcommand{\MetricBangBangBfive}{0.2212}
\newcommand{\MetricBangBangBseven}{0.2360}
\newcommand{\MetricBangBangBeight}{0.2484}
\newcommand{\MetricBangBangBgena}{\textbf{??}}
\newcommand{\MetricBangBangBgenb}{\textbf{??}}
\newcommand{\DecompBangBangComfort}{0.1148}
\newcommand{\DecompBangBangComfortFairbanks}{0.0840}
\newcommand{\DecompBangBangComfortDenver}{0.0995}
\newcommand{\DecompBangBangComfortAtlanta}{0.1210}
\newcommand{\DecompBangBangComfortPhoenix}{0.1268}
\newcommand{\DecompBangBangComfortSingapore}{0.1426}
\newcommand{\DecompBangBangEnergy}{0.0607}
\newcommand{\DecompBangBangEnergyFairbanks}{0.0520}
\newcommand{\DecompBangBangEnergyDenver}{0.0451}
\newcommand{\DecompBangBangEnergyAtlanta}{0.0612}
\newcommand{\DecompBangBangEnergyPhoenix}{0.0682}
\newcommand{\DecompBangBangEnergySingapore}{0.0769}
\newcommand{\DecompBangBangCarbon}{0.0491}
\newcommand{\DecompBangBangCarbonFairbanks}{0.0636}
\newcommand{\DecompBangBangCarbonDenver}{0.0384}
\newcommand{\DecompBangBangCarbonAtlanta}{0.0494}
\newcommand{\DecompBangBangCarbonPhoenix}{0.0418}
\newcommand{\DecompBangBangCarbonSingapore}{0.0523}
\newcommand{\SecPerStepBangBang}{0.056}
\newcommand{\PolicyMsBangBang}{0.13}
\newcommand{\MetricPidWaterOff}{0.0541}
\newcommand{\MetricSdPidWaterOff}{0.0448}
\newcommand{\MetricPidWaterOffFairbanks}{0.0785}
\newcommand{\MetricPidWaterOffDenver}{0.0374}
\newcommand{\MetricPidWaterOffAtlanta}{0.0301}
\newcommand{\MetricPidWaterOffPhoenix}{0.0781}
\newcommand{\MetricPidWaterOffSingapore}{0.0466}
\newcommand{\MetricPidWaterOffBone}{0.0855}
\newcommand{\MetricPidWaterOffBtwo}{0.0472}
\newcommand{\MetricPidWaterOffBfour}{0.0352}
\newcommand{\MetricPidWaterOffBfive}{0.0459}
\newcommand{\MetricPidWaterOffBseven}{0.0427}
\newcommand{\MetricPidWaterOffBeight}{0.0684}
\newcommand{\MetricPidWaterOffBgena}{\textbf{??}}
\newcommand{\MetricPidWaterOffBgenb}{\textbf{??}}
\newcommand{\DecompPidWaterOffComfort}{\textbf{??}}
\newcommand{\DecompPidWaterOffComfortFairbanks}{\textbf{??}}
\newcommand{\DecompPidWaterOffComfortDenver}{\textbf{??}}
\newcommand{\DecompPidWaterOffComfortAtlanta}{\textbf{??}}
\newcommand{\DecompPidWaterOffComfortPhoenix}{\textbf{??}}
\newcommand{\DecompPidWaterOffComfortSingapore}{\textbf{??}}
\newcommand{\DecompPidWaterOffEnergy}{\textbf{??}}
\newcommand{\DecompPidWaterOffEnergyFairbanks}{\textbf{??}}
\newcommand{\DecompPidWaterOffEnergyDenver}{\textbf{??}}
\newcommand{\DecompPidWaterOffEnergyAtlanta}{\textbf{??}}
\newcommand{\DecompPidWaterOffEnergyPhoenix}{\textbf{??}}
\newcommand{\DecompPidWaterOffEnergySingapore}{\textbf{??}}
\newcommand{\DecompPidWaterOffCarbon}{\textbf{??}}
\newcommand{\DecompPidWaterOffCarbonFairbanks}{\textbf{??}}
\newcommand{\DecompPidWaterOffCarbonDenver}{\textbf{??}}
\newcommand{\DecompPidWaterOffCarbonAtlanta}{\textbf{??}}
\newcommand{\DecompPidWaterOffCarbonPhoenix}{\textbf{??}}
\newcommand{\DecompPidWaterOffCarbonSingapore}{\textbf{??}}
\newcommand{\SecPerStepPidWaterOff}{\textbf{??}}
\newcommand{\PolicyMsPidWaterOff}{\textbf{??}}
\newcommand{\MetricPidOff}{0.0674}
\newcommand{\MetricSdPidOff}{0.0476}
\newcommand{\MetricPidOffFairbanks}{0.0755}
\newcommand{\MetricPidOffDenver}{0.0496}
\newcommand{\MetricPidOffAtlanta}{0.0418}
\newcommand{\MetricPidOffPhoenix}{0.0950}
\newcommand{\MetricPidOffSingapore}{0.0750}
\newcommand{\MetricPidOffBone}{0.0786}
\newcommand{\MetricPidOffBtwo}{0.0662}
\newcommand{\MetricPidOffBfour}{0.0394}
\newcommand{\MetricPidOffBfive}{0.0716}
\newcommand{\MetricPidOffBseven}{0.0582}
\newcommand{\MetricPidOffBeight}{0.0899}
\newcommand{\MetricPidOffBgena}{\textbf{??}}
\newcommand{\MetricPidOffBgenb}{\textbf{??}}
\newcommand{\DecompPidOffComfort}{\textbf{??}}
\newcommand{\DecompPidOffComfortFairbanks}{\textbf{??}}
\newcommand{\DecompPidOffComfortDenver}{\textbf{??}}
\newcommand{\DecompPidOffComfortAtlanta}{\textbf{??}}
\newcommand{\DecompPidOffComfortPhoenix}{\textbf{??}}
\newcommand{\DecompPidOffComfortSingapore}{\textbf{??}}
\newcommand{\DecompPidOffEnergy}{\textbf{??}}
\newcommand{\DecompPidOffEnergyFairbanks}{\textbf{??}}
\newcommand{\DecompPidOffEnergyDenver}{\textbf{??}}
\newcommand{\DecompPidOffEnergyAtlanta}{\textbf{??}}
\newcommand{\DecompPidOffEnergyPhoenix}{\textbf{??}}
\newcommand{\DecompPidOffEnergySingapore}{\textbf{??}}
\newcommand{\DecompPidOffCarbon}{\textbf{??}}
\newcommand{\DecompPidOffCarbonFairbanks}{\textbf{??}}
\newcommand{\DecompPidOffCarbonDenver}{\textbf{??}}
\newcommand{\DecompPidOffCarbonAtlanta}{\textbf{??}}
\newcommand{\DecompPidOffCarbonPhoenix}{\textbf{??}}
\newcommand{\DecompPidOffCarbonSingapore}{\textbf{??}}
\newcommand{\SecPerStepPidOff}{\textbf{??}}
\newcommand{\PolicyMsPidOff}{\textbf{??}}
\newcommand{\MetricFixedSpOff}{0.0659}
\newcommand{\MetricSdFixedSpOff}{0.0540}
\newcommand{\MetricFixedSpOffFairbanks}{0.0743}
\newcommand{\MetricFixedSpOffDenver}{0.0768}
\newcommand{\MetricFixedSpOffAtlanta}{0.0331}
\newcommand{\MetricFixedSpOffPhoenix}{0.1202}
\newcommand{\MetricFixedSpOffSingapore}{0.0252}
\newcommand{\MetricFixedSpOffBone}{0.0721}
\newcommand{\MetricFixedSpOffBtwo}{0.0804}
\newcommand{\MetricFixedSpOffBfour}{0.0349}
\newcommand{\MetricFixedSpOffBfive}{0.0430}
\newcommand{\MetricFixedSpOffBseven}{0.0919}
\newcommand{\MetricFixedSpOffBeight}{0.0731}
\newcommand{\MetricFixedSpOffBgena}{\textbf{??}}
\newcommand{\MetricFixedSpOffBgenb}{\textbf{??}}
\newcommand{\DecompFixedSpOffComfort}{\textbf{??}}
\newcommand{\DecompFixedSpOffComfortFairbanks}{\textbf{??}}
\newcommand{\DecompFixedSpOffComfortDenver}{\textbf{??}}
\newcommand{\DecompFixedSpOffComfortAtlanta}{\textbf{??}}
\newcommand{\DecompFixedSpOffComfortPhoenix}{\textbf{??}}
\newcommand{\DecompFixedSpOffComfortSingapore}{\textbf{??}}
\newcommand{\DecompFixedSpOffEnergy}{\textbf{??}}
\newcommand{\DecompFixedSpOffEnergyFairbanks}{\textbf{??}}
\newcommand{\DecompFixedSpOffEnergyDenver}{\textbf{??}}
\newcommand{\DecompFixedSpOffEnergyAtlanta}{\textbf{??}}
\newcommand{\DecompFixedSpOffEnergyPhoenix}{\textbf{??}}
\newcommand{\DecompFixedSpOffEnergySingapore}{\textbf{??}}
\newcommand{\DecompFixedSpOffCarbon}{\textbf{??}}
\newcommand{\DecompFixedSpOffCarbonFairbanks}{\textbf{??}}
\newcommand{\DecompFixedSpOffCarbonDenver}{\textbf{??}}
\newcommand{\DecompFixedSpOffCarbonAtlanta}{\textbf{??}}
\newcommand{\DecompFixedSpOffCarbonPhoenix}{\textbf{??}}
\newcommand{\DecompFixedSpOffCarbonSingapore}{\textbf{??}}
\newcommand{\SecPerStepFixedSpOff}{\textbf{??}}
\newcommand{\PolicyMsFixedSpOff}{\textbf{??}}
\newcommand{\MetricBangBangOff}{0.2252}
\newcommand{\MetricSdBangBangOff}{0.0586}
\newcommand{\MetricBangBangOffFairbanks}{0.1999}
\newcommand{\MetricBangBangOffDenver}{0.1844}
\newcommand{\MetricBangBangOffAtlanta}{0.2320}
\newcommand{\MetricBangBangOffPhoenix}{0.2378}
\newcommand{\MetricBangBangOffSingapore}{0.2719}
\newcommand{\MetricBangBangOffBone}{0.2114}
\newcommand{\MetricBangBangOffBtwo}{0.2432}
\newcommand{\MetricBangBangOffBfour}{0.1885}
\newcommand{\MetricBangBangOffBfive}{0.2232}
\newcommand{\MetricBangBangOffBseven}{0.2354}
\newcommand{\MetricBangBangOffBeight}{0.2495}
\newcommand{\MetricBangBangOffBgena}{\textbf{??}}
\newcommand{\MetricBangBangOffBgenb}{\textbf{??}}
\newcommand{\DecompBangBangOffComfort}{\textbf{??}}
\newcommand{\DecompBangBangOffComfortFairbanks}{\textbf{??}}
\newcommand{\DecompBangBangOffComfortDenver}{\textbf{??}}
\newcommand{\DecompBangBangOffComfortAtlanta}{\textbf{??}}
\newcommand{\DecompBangBangOffComfortPhoenix}{\textbf{??}}
\newcommand{\DecompBangBangOffComfortSingapore}{\textbf{??}}
\newcommand{\DecompBangBangOffEnergy}{\textbf{??}}
\newcommand{\DecompBangBangOffEnergyFairbanks}{\textbf{??}}
\newcommand{\DecompBangBangOffEnergyDenver}{\textbf{??}}
\newcommand{\DecompBangBangOffEnergyAtlanta}{\textbf{??}}
\newcommand{\DecompBangBangOffEnergyPhoenix}{\textbf{??}}
\newcommand{\DecompBangBangOffEnergySingapore}{\textbf{??}}
\newcommand{\DecompBangBangOffCarbon}{\textbf{??}}
\newcommand{\DecompBangBangOffCarbonFairbanks}{\textbf{??}}
\newcommand{\DecompBangBangOffCarbonDenver}{\textbf{??}}
\newcommand{\DecompBangBangOffCarbonAtlanta}{\textbf{??}}
\newcommand{\DecompBangBangOffCarbonPhoenix}{\textbf{??}}
\newcommand{\DecompBangBangOffCarbonSingapore}{\textbf{??}}
\newcommand{\SecPerStepBangBangOff}{\textbf{??}}
\newcommand{\PolicyMsBangBangOff}{\textbf{??}}
\newcommand{\DiffOurs}{\textbf{??}}
\newcommand{\DiffLoOurs}{\textbf{??}}
\newcommand{\DiffHiOurs}{\textbf{??}}
\newcommand{\RatioOurs}{\textbf{??}}
\newcommand{\RatioLoOurs}{\textbf{??}}
\newcommand{\RatioHiOurs}{\textbf{??}}
\newcommand{\RelRedOurs}{\textbf{??}}
\newcommand{\RelRedLoOurs}{\textbf{??}}
\newcommand{\RelRedHiOurs}{\textbf{??}}
\newcommand{\WinsOurs}{\textbf{??}}
\newcommand{\LossesOurs}{\textbf{??}}
\newcommand{\PWilcoxonOurs}{\textbf{??}}
\newcommand{\DiffOursFairbanks}{\textbf{??}}
\newcommand{\DiffOursDenver}{\textbf{??}}
\newcommand{\DiffOursAtlanta}{\textbf{??}}
\newcommand{\DiffOursPhoenix}{\textbf{??}}
\newcommand{\DiffOursSingapore}{\textbf{??}}
\newcommand{\DiffOursVFour}{-0.0009}
\newcommand{\DiffLoOursVFour}{-0.0014}
\newcommand{\DiffHiOursVFour}{-0.0004}
\newcommand{\RatioOursVFour}{0.903}
\newcommand{\RatioLoOursVFour}{0.848}
\newcommand{\RatioHiOursVFour}{0.949}
\newcommand{\RelRedOursVFour}{9.7\%}
\newcommand{\RelRedLoOursVFour}{5.1\%}
\newcommand{\RelRedHiOursVFour}{15.2\%}
\newcommand{\WinsOursVFour}{19}
\newcommand{\LossesOursVFour}{11}
\newcommand{\PWilcoxonOursVFour}{<0.01}
\newcommand{\DiffOursVFourFairbanks}{-0.0003}
\newcommand{\DiffOursVFourDenver}{-0.0000}
\newcommand{\DiffOursVFourAtlanta}{-0.0003}
\newcommand{\DiffOursVFourPhoenix}{-0.0020}
\newcommand{\DiffOursVFourSingapore}{-0.0018}
\newcommand{\DiffPidWater}{-0.0467}
\newcommand{\DiffLoPidWater}{-0.0616}
\newcommand{\DiffHiPidWater}{-0.0342}
\newcommand{\RatioPidWater}{0.151}
\newcommand{\RatioLoPidWater}{0.117}
\newcommand{\RatioHiPidWater}{0.188}
\newcommand{\RelRedPidWater}{84.9\%}
\newcommand{\RelRedLoPidWater}{81.2\%}
\newcommand{\RelRedHiPidWater}{88.3\%}
\newcommand{\WinsPidWater}{30}
\newcommand{\LossesPidWater}{0}
\newcommand{\PWilcoxonPidWater}{<10^{-4}}
\newcommand{\DiffPidWaterFairbanks}{-0.0616}
\newcommand{\DiffPidWaterDenver}{-0.0330}
\newcommand{\DiffPidWaterAtlanta}{-0.0252}
\newcommand{\DiffPidWaterPhoenix}{-0.0725}
\newcommand{\DiffPidWaterSingapore}{-0.0411}
\newcommand{\DiffPid}{-0.0595}
\newcommand{\DiffLoPid}{-0.0752}
\newcommand{\DiffHiPid}{-0.0453}
\newcommand{\RatioPid}{0.123}
\newcommand{\RatioLoPid}{0.088}
\newcommand{\RatioHiPid}{0.169}
\newcommand{\RelRedPid}{87.7\%}
\newcommand{\RelRedLoPid}{83.1\%}
\newcommand{\RelRedHiPid}{91.2\%}
\newcommand{\WinsPid}{30}
\newcommand{\LossesPid}{0}
\newcommand{\PWilcoxonPid}{<10^{-4}}
\newcommand{\DiffPidFairbanks}{-0.0584}
\newcommand{\DiffPidDenver}{-0.0457}
\newcommand{\DiffPidAtlanta}{-0.0377}
\newcommand{\DiffPidPhoenix}{-0.0879}
\newcommand{\DiffPidSingapore}{-0.0679}
\newcommand{\DiffFixedSp}{-0.0572}
\newcommand{\DiffLoFixedSp}{-0.0748}
\newcommand{\DiffHiFixedSp}{-0.0402}
\newcommand{\RatioFixedSp}{0.127}
\newcommand{\RatioLoFixedSp}{0.088}
\newcommand{\RatioHiFixedSp}{0.179}
\newcommand{\RelRedFixedSp}{87.3\%}
\newcommand{\RelRedLoFixedSp}{82.1\%}
\newcommand{\RelRedHiFixedSp}{91.2\%}
\newcommand{\WinsFixedSp}{30}
\newcommand{\LossesFixedSp}{0}
\newcommand{\PWilcoxonFixedSp}{<10^{-4}}
\newcommand{\DiffFixedSpFairbanks}{-0.0583}
\newcommand{\DiffFixedSpDenver}{-0.0722}
\newcommand{\DiffFixedSpAtlanta}{-0.0275}
\newcommand{\DiffFixedSpPhoenix}{-0.1086}
\newcommand{\DiffFixedSpSingapore}{-0.0193}
\newcommand{\DiffBangBang}{-0.2162}
\newcommand{\DiffLoBangBang}{-0.2356}
\newcommand{\DiffHiBangBang}{-0.1959}
\newcommand{\RatioBangBang}{0.037}
\newcommand{\RatioLoBangBang}{0.025}
\newcommand{\RatioHiBangBang}{0.053}
\newcommand{\RelRedBangBang}{96.3\%}
\newcommand{\RelRedLoBangBang}{94.7\%}
\newcommand{\RelRedHiBangBang}{97.5\%}
\newcommand{\WinsBangBang}{30}
\newcommand{\LossesBangBang}{0}
\newcommand{\PWilcoxonBangBang}{<10^{-4}}
\newcommand{\DiffBangBangFairbanks}{-0.1821}
\newcommand{\DiffBangBangDenver}{-0.1780}
\newcommand{\DiffBangBangAtlanta}{-0.2260}
\newcommand{\DiffBangBangPhoenix}{-0.2292}
\newcommand{\DiffBangBangSingapore}{-0.2658}
\newcommand{\DiffPidWaterOff}{-0.0458}
\newcommand{\DiffLoPidWaterOff}{-0.0606}
\newcommand{\DiffHiPidWaterOff}{-0.0336}
\newcommand{\RatioPidWaterOff}{0.153}
\newcommand{\RatioLoPidWaterOff}{0.119}
\newcommand{\RatioHiPidWaterOff}{0.191}
\newcommand{\RelRedPidWaterOff}{84.7\%}
\newcommand{\RelRedLoPidWaterOff}{80.9\%}
\newcommand{\RelRedHiPidWaterOff}{88.1\%}
\newcommand{\WinsPidWaterOff}{30}
\newcommand{\LossesPidWaterOff}{0}
\newcommand{\PWilcoxonPidWaterOff}{<10^{-4}}
\newcommand{\DiffPidWaterOffFairbanks}{-0.0609}
\newcommand{\DiffPidWaterOffDenver}{-0.0324}
\newcommand{\DiffPidWaterOffAtlanta}{-0.0246}
\newcommand{\DiffPidWaterOffPhoenix}{-0.0706}
\newcommand{\DiffPidWaterOffSingapore}{-0.0406}
\newcommand{\DiffPidOff}{-0.0590}
\newcommand{\DiffLoPidOff}{-0.0747}
\newcommand{\DiffHiPidOff}{-0.0448}
\newcommand{\RatioPidOff}{0.123}
\newcommand{\RatioLoPidOff}{0.088}
\newcommand{\RatioHiPidOff}{0.170}
\newcommand{\RelRedPidOff}{87.7\%}
\newcommand{\RelRedLoPidOff}{83.0\%}
\newcommand{\RelRedHiPidOff}{91.2\%}
\newcommand{\WinsPidOff}{30}
\newcommand{\LossesPidOff}{0}
\newcommand{\PWilcoxonPidOff}{<10^{-4}}
\newcommand{\DiffPidOffFairbanks}{-0.0579}
\newcommand{\DiffPidOffDenver}{-0.0446}
\newcommand{\DiffPidOffAtlanta}{-0.0363}
\newcommand{\DiffPidOffPhoenix}{-0.0874}
\newcommand{\DiffPidOffSingapore}{-0.0689}
\newcommand{\DiffFixedSpOff}{-0.0576}
\newcommand{\DiffLoFixedSpOff}{-0.0760}
\newcommand{\DiffHiFixedSpOff}{-0.0403}
\newcommand{\RatioFixedSpOff}{0.126}
\newcommand{\RatioLoFixedSpOff}{0.087}
\newcommand{\RatioHiFixedSpOff}{0.180}
\newcommand{\RelRedFixedSpOff}{87.4\%}
\newcommand{\RelRedLoFixedSpOff}{82.0\%}
\newcommand{\RelRedHiFixedSpOff}{91.3\%}
\newcommand{\WinsFixedSpOff}{30}
\newcommand{\LossesFixedSpOff}{0}
\newcommand{\PWilcoxonFixedSpOff}{<10^{-4}}
\newcommand{\DiffFixedSpOffFairbanks}{-0.0567}
\newcommand{\DiffFixedSpOffDenver}{-0.0719}
\newcommand{\DiffFixedSpOffAtlanta}{-0.0276}
\newcommand{\DiffFixedSpOffPhoenix}{-0.1127}
\newcommand{\DiffFixedSpOffSingapore}{-0.0192}
\newcommand{\DiffBangBangOff}{-0.2169}
\newcommand{\DiffLoBangBangOff}{-0.2361}
\newcommand{\DiffHiBangBangOff}{-0.1969}
\newcommand{\RatioBangBangOff}{0.037}
\newcommand{\RatioLoBangBangOff}{0.025}
\newcommand{\RatioHiBangBangOff}{0.053}
\newcommand{\RelRedBangBangOff}{96.3\%}
\newcommand{\RelRedLoBangBangOff}{94.7\%}
\newcommand{\RelRedHiBangBangOff}{97.5\%}
\newcommand{\WinsBangBangOff}{30}
\newcommand{\LossesBangBangOff}{0}
\newcommand{\PWilcoxonBangBangOff}{<10^{-4}}
\newcommand{\DiffBangBangOffFairbanks}{-0.1824}
\newcommand{\DiffBangBangOffDenver}{-0.1794}
\newcommand{\DiffBangBangOffAtlanta}{-0.2265}
\newcommand{\DiffBangBangOffPhoenix}{-0.2303}
\newcommand{\DiffBangBangOffSingapore}{-0.2659}
\newcommand{\NBoot}{5000}
\newcommand{\GapMeanOurs}{\textbf{??}}
\newcommand{\GapMinOurs}{\textbf{??}}
\newcommand{\GapMaxOurs}{\textbf{??}}
\newcommand{\GapMeanOursVFour}{\textbf{??}}
\newcommand{\GapMinOursVFour}{\textbf{??}}
\newcommand{\GapMaxOursVFour}{\textbf{??}}
\newcommand{\GapMeanPidWater}{-1.7\%}
\newcommand{\GapMinPidWater}{-12.7\%}
\newcommand{\GapMaxPidWater}{+4.6\%}
\newcommand{\GapMeanPid}{-1.3\%}
\newcommand{\GapMinPid}{-9.6\%}
\newcommand{\GapMaxPid}{+7.7\%}
\newcommand{\GapMeanFixedSp}{-0.4\%}
\newcommand{\GapMinFixedSp}{-4.4\%}
\newcommand{\GapMaxFixedSp}{+16.1\%}
\newcommand{\GapMeanBangBang}{+0.5\%}
\newcommand{\GapMinBangBang}{-2.4\%}
\newcommand{\GapMaxBangBang}{+5.7\%}
\newcommand{\GapMeanPidWaterOff}{\textbf{??}}
\newcommand{\GapMinPidWaterOff}{\textbf{??}}
\newcommand{\GapMaxPidWaterOff}{\textbf{??}}
\newcommand{\GapMeanPidOff}{\textbf{??}}
\newcommand{\GapMinPidOff}{\textbf{??}}
\newcommand{\GapMaxPidOff}{\textbf{??}}
\newcommand{\GapMeanFixedSpOff}{\textbf{??}}
\newcommand{\GapMinFixedSpOff}{\textbf{??}}
\newcommand{\GapMaxFixedSpOff}{\textbf{??}}
\newcommand{\GapMeanBangBangOff}{\textbf{??}}
\newcommand{\GapMinBangBangOff}{\textbf{??}}
\newcommand{\GapMaxBangBangOff}{\textbf{??}}
\newcommand{\WorstScen}{building\_train\_5/fairbanks\_5}
\newcommand{\WorstScenDiff}{-0.0145}
\newcommand{\BestScen}{building\_train\_1/fairbanks\_1}
\newcommand{\BestScenDiff}{-0.1635}
\newcommand{\NScenWorse}{0}
\newcommand{\MaxOursScore}{0.0506}
\newcommand{\MaxOursScen}{building\_train\_1/fairbanks\_1}
\newcommand{\MaxBaseScore}{0.2141}
\newcommand{\ZonesBone}{39}
\newcommand{\AhusBone}{4}
\newcommand{\ZonesBtwo}{85}
\newcommand{\AhusBtwo}{9}
\newcommand{\ZonesBfour}{54}
\newcommand{\AhusBfour}{6}
\newcommand{\ZonesBfive}{79}
\newcommand{\AhusBfive}{8}
\newcommand{\ZonesBseven}{36}
\newcommand{\AhusBseven}{4}
\newcommand{\ZonesBeight}{30}
\newcommand{\AhusBeight}{3}
\newcommand{\ZonesBgena}{29}
\newcommand{\AhusBgena}{5}
\newcommand{\ZonesBgenb}{28}
\newcommand{\AhusBgenb}{4}
\newcommand{\NHeld}{10}
\newcommand{\NHeldBuild}{2}
\newcommand{\HeldMetricOursBgena}{0.0044}
\newcommand{\HeldMetricOursBgenb}{0.0177}
\newcommand{\HeldMetricOursFairbanks}{0.0157}
\newcommand{\HeldMetricOursDenver}{0.0184}
\newcommand{\HeldMetricOursAtlanta}{0.0071}
\newcommand{\HeldMetricOursPhoenix}{0.0056}
\newcommand{\HeldMetricOursSingapore}{0.0086}
\newcommand{\HeldMetricPidWaterBgena}{0.0310}
\newcommand{\HeldMetricPidWaterBgenb}{0.0488}
\newcommand{\HeldMetricPidWaterFairbanks}{0.0515}
\newcommand{\HeldMetricPidWaterDenver}{0.0488}
\newcommand{\HeldMetricPidWaterAtlanta}{0.0316}
\newcommand{\HeldMetricPidWaterPhoenix}{0.0273}
\newcommand{\HeldMetricPidWaterSingapore}{0.0402}
\newcommand{\HeldMetricOursVFourBgena}{0.0047}
\newcommand{\HeldMetricOursVFourBgenb}{0.0181}
\newcommand{\HeldMetricOursVFourFairbanks}{0.0154}
\newcommand{\HeldMetricOursVFourDenver}{0.0185}
\newcommand{\HeldMetricOursVFourAtlanta}{0.0075}
\newcommand{\HeldMetricOursVFourPhoenix}{0.0057}
\newcommand{\HeldMetricOursVFourSingapore}{0.0099}
\newcommand{\HeldMetricOurs}{0.0111}
\newcommand{\HeldMetricPidWater}{0.0399}
\newcommand{\HeldMetricOursVFour}{0.0114}
\newcommand{\HeldDiff}{-0.0288}
\newcommand{\HeldDiffLo}{-0.0331}
\newcommand{\HeldDiffHi}{-0.0248}
\newcommand{\HeldRatio}{0.278}
\newcommand{\HeldRatioLo}{0.175}
\newcommand{\HeldRatioHi}{0.355}
\newcommand{\HeldRelRed}{72.2\%}
\newcommand{\HeldWins}{10}
\newcommand{\HeldPWilcoxon}{<0.01}
\newcommand{\HeldVFourDiff}{-0.00030}
\newcommand{\HeldVFourDiffLo}{-0.00071}
\newcommand{\HeldVFourDiffHi}{+0.00006}
\newcommand{\HeldVFourRatio}{0.974}
\newcommand{\HeldVFourWins}{6}
\newcommand{\HeldDecompOursComfort}{0.0018}
\newcommand{\HeldDecompPidWaterComfort}{0.0267}
\newcommand{\HeldDecompOursEnergy}{0.0045}
\newcommand{\HeldDecompPidWaterEnergy}{0.0074}
\newcommand{\HeldDecompOursCarbon}{0.0048}
\newcommand{\HeldDecompPidWaterCarbon}{0.0058}
\newcommand{\NTuneConfigs}{16}
\newcommand{\NTuneRandom}{15}
\newcommand{\NAblation}{8}
\newcommand{\AblRevKi}{0.00905}
\newcommand{\AblRelRevKi}{+6.7\%}
\newcommand{\AblRelDefRevKi}{-4.4\%}
\newcommand{\AblAddKi}{0.00855}
\newcommand{\AblRelAddKi}{+0.7\%}
\newcommand{\AblRelDefAddKi}{-9.8\%}
\newcommand{\AblRevMarginPre}{0.00846}
\newcommand{\AblRelRevMarginPre}{-0.3\%}
\newcommand{\AblRelDefRevMarginPre}{-10.7\%}
\newcommand{\AblAddMarginPre}{0.00910}
\newcommand{\AblRelAddMarginPre}{+7.2\%}
\newcommand{\AblRelDefAddMarginPre}{-4.0\%}
\newcommand{\AblRevSatMix}{0.00849}
\newcommand{\AblRelRevSatMix}{+0.0\%}
\newcommand{\AblRelDefRevSatMix}{-10.4\%}
\newcommand{\AblAddSatMix}{0.00947}
\newcommand{\AblRelAddSatMix}{+11.6\%}
\newcommand{\AblRelDefAddSatMix}{+0.0\%}
\newcommand{\AblRevSbHyst}{0.00850}
\newcommand{\AblRelRevSbHyst}{+0.2\%}
\newcommand{\AblRelDefRevSbHyst}{-10.2\%}
\newcommand{\AblAddSbHyst}{0.00944}
\newcommand{\AblRelAddSbHyst}{+11.2\%}
\newcommand{\AblRelDefAddSbHyst}{-0.3\%}
\newcommand{\NCandV}{13}
\newcommand{\CandBestScreen}{0.00758}
\newcommand{\CandSatColdScreen}{0.00770}
\newcommand{\CandNightEightScreen}{0.00846}
\newcommand{\TuneBestScreen}{0.0085}
\newcommand{\TuneWorstScreen}{0.0120}
\newcommand{\TuneDefaultScreen}{0.0095}
\newcommand{\TunePidWaterScreen}{0.0490}
\newcommand{\ParSource}{tuned}
\newcommand{\ParMarginOcc}{0.3}
\newcommand{\ParMarginPre}{0.8}
\newcommand{\ParNightLowExt}{8}
\newcommand{\ParNightHighExt}{6}
\newcommand{\ParKqWPerK}{400}
\newcommand{\ParKiWPerKStep}{120}
\newcommand{\ParSatCoolHi}{285}
\newcommand{\ParSatMixFloor}{292}
\newcommand{\ParLeadSafety}{1.3}
\newcommand{\ParRateZeroKph}{2}
\newcommand{\ParHolidayCheckH}{7}
\newcommand{\ParHeatSeasonSteps}{288}
\newcommand{\ParMinOnSteps}{3}
\newcommand{\ParDischMarginK}{1}
\newcommand{\ParRateEma}{0.5}
\newcommand{\ParLeadMinH}{0.75}
\newcommand{\ParLeadMaxH}{8}
\newcommand{\ParValveTarget}{0.9}
\newcommand{\ParQCapW}{12000}
\newcommand{\ParSbHystK}{1}
\newcommand{\ParQSbW}{4000}
\newcommand{\VthreeAhuOnUnocc}{84.9\%}
\newcommand{\PhoenixPeakOatC}{46}
\newcommand{\ExtremeRatioFairbanks}{2.1}
\newcommand{\ExtremeRatioPhoenix}{2.1}
\newcommand{\TankLossExtraKw}{0.25}
\newcommand{\NAnalysis}{30}
\newcommand{\UnoccShareOurs}{25.4\%}
\newcommand{\RegretShareOurs}{4.2\%}
\newcommand{\EnergyShareOurs}{56.4\%}
\newcommand{\CarbonShareOurs}{39.3\%}
\newcommand{\OccOutOurs}{3.2\%}
\newcommand{\OccOutGtOneOurs}{0.1\%}
\newcommand{\OccWorstDevOurs}{0.11}
\newcommand{\FirstWorkdayRatioOurs}{1.03}
\newcommand{\FirstWorkdayRatioMedianOurs}{0.65}
\newcommand{\UnoccShareOursFairbanks}{37.7\%}
\newcommand{\UnoccShareOursDenver}{27.9\%}
\newcommand{\UnoccShareOursAtlanta}{24.8\%}
\newcommand{\UnoccShareOursPhoenix}{20.6\%}
\newcommand{\UnoccShareOursSingapore}{16.0\%}
\newcommand{\PeakLossHourOurs}{5}
\newcommand{\PeakLossShareOurs}{12.5\%}
\newcommand{\UnoccSharePidWater}{5.1\%}
\newcommand{\RegretSharePidWater}{71.5\%}
\newcommand{\EnergySharePidWater}{17.9\%}
\newcommand{\CarbonSharePidWater}{10.6\%}
\newcommand{\OccOutPidWater}{61.5\%}
\newcommand{\OccOutGtOnePidWater}{25.6\%}
\newcommand{\OccWorstDevPidWater}{3.21}
\newcommand{\FirstWorkdayRatioPidWater}{0.80}
\newcommand{\FirstWorkdayRatioMedianPidWater}{0.81}
\newcommand{\UnoccSharePidWaterFairbanks}{11.5\%}
\newcommand{\UnoccSharePidWaterDenver}{1.6\%}
\newcommand{\UnoccSharePidWaterAtlanta}{3.4\%}
\newcommand{\UnoccSharePidWaterPhoenix}{3.5\%}
\newcommand{\UnoccSharePidWaterSingapore}{5.7\%}
\newcommand{\PeakLossHourPidWater}{6}
\newcommand{\PeakLossSharePidWater}{9.3\%}
\newcommand{\NFirstWorkdayWorse}{13}
\newcommand{\NFirstWorkdayTotal}{30}
\newcommand{\LbDate}{2026-10-22}
\newcommand{\LbNTeams}{292}
\newcommand{\LbRankOne}{0.00578}
\newcommand{\LbRankEight}{0.00647}
\newcommand{\LbRankTen}{0.00658}
\newcommand{\LbRankThirty}{0.00996}
\newcommand{\LbRankFortyNine}{0.01056}
\newcommand{\OursVsLbOne}{1.4}
\newcommand{\NCandFull}{13}
\newcommand{\CandWinner}{c10\_satcold\_night8\_hi6}
\newcommand{\CandWinnerMetric}{0.00831}
\newcommand{\CandWinnerVFourMetric}{0.00920}
\newcommand{\CandWinnerRatio}{0.903}
\newcommand{\CandWinnerDiff}{-0.00089}
\newcommand{\CandWinnerCiLo}{-0.00139}
\newcommand{\CandWinnerCiHi}{-0.00044}
\newcommand{\CandWinnerWins}{19}
\newcommand{\CandSatColdFull}{0.00856}
\newcommand{\CandNightEightFull}{0.00908}
\newcommand{\VFourSatCoolHi}{293}
\newcommand{\VFourNightLowExt}{4}
\newcommand{\VFourNightHighExt}{2}
\newcommand{\VFourMarginOcc}{0.3}
\newcommand{\LocalVFourMetric}{0.0092}
\newcommand{\LocalPidWaterMetric}{0.0550}
\newcommand{\MachineGapMean}{0.000\%}
\newcommand{\MachineGapMax}{0.000\%}
\newcommand{\MachineGapN}{30}
\newcommand{\MachineGapPidWater}{0.017\%}
\newcommand{\SimSecPerStep}{0.055}
\newcommand{\SimMinPerScenario}{8.0}
\newcommand{\TotalCoreHours}{134.9}

% Figure include that degrades to a visible placeholder when the file is not built yet.
\newcommand{\figfile}[2]{\IfFileExists{#2}{\includegraphics[width=#1]{#2}}{\fbox{\parbox{#1}{\centering\textbf{??} missing figure \texttt{\detokenize{#2}}}}}}

\title{Schedule-Aware Demand-Based Reset: A Structured Rule-Based Controller for Cross-Building HVAC Transfer in the Smart Buildings Challenge}

\author{%
  {[Author Name]}\\
  {[Affiliation]}\\
  \texttt{neelsirivara@outlook.com}
}

\begin{document}

\maketitle

\begin{abstract}
The NeurIPS 2026 Smart Buildings Challenge asks for one HVAC control policy that operates
procedurally generated commercial buildings it has never seen, under 30-day weather strips
from five climates, trading occupant comfort against energy cost and carbon. We present
SADR (Schedule-Aware Demand-based Reset), a pure-numpy rule-based controller that reads the
plant topology from the building metadata and is otherwise built from facts shared by every
generated building: comfort is charged only while zones are occupied, fan power grows as
static pressure to the three-halves, the chiller has a constant coefficient of performance,
and all heat reaches a zone through a reheat coil whose discharge is capped at 308.15\,K.
SADR gates its mode on the occupancy count and schedule flags (with a self-tuning optimum
start), regulates each zone with a PI loop on \emph{net heat} translated into damper and
valve commands through the supply-air and discharge physics, and resets static pressure,
loop differential pressure and hot-water temperature from zone demand. On the 30 public
scenarios SADR scores \MetricOurs{} on the leaderboard metric (mean absolute per-step reward,
lower is better) against \MetricPidWater{} for the strongest organizer baseline
(\texttt{pid\_water}) under identical conditions, a paired-bootstrap ratio of
\RatioPidWater{} (95\% CI \RatioLoPidWater--\RatioHiPidWater; \WinsPidWater{}/\NScen{} wins); on
\NHeld{} scenarios from \NHeldBuild{} freshly generated buildings and unseen strips the ratio
is \HeldRatio. Decomposing the remaining loss by term, station and hour led to one
physics-motivated revision (coldest supply air whenever a handler cools, deeper night
setback) that cut the public-set loss by a further \RelRedOursVFour{} (a smaller, not
significant gain on the held-out set). We release the controller and report where it still
loses.
\end{abstract}

\section{Introduction}

Learned building controllers have been studied for two decades, yet few run in real
buildings; the recurring obstacles are data-hungry training, fragile robustness and poor
generalization across buildings~\citep{wang2020rl}. The Smart Buildings Challenge (SB26)
isolates the last of these: a submission is a single policy class, instantiated on held-out
buildings drawn from the same procedural generator as six public ones and scored on 30-day
episodes under weather strips it has not seen~\citep{sb26kit}. The observation and action
vectors change shape with the building ($K$ air handlers and $V$ zones expose $K{+}1$ on/off
commands and $3K{+}2{+}2V$ continuous setpoints), and the simulator is the open-source Smart
Buildings simulator behind Google's calibrated digital twins and the Smart Buildings Control
Suite~\citep{goldfeder2023lightweight,goldfeder2025suite}.

We take the position that the right first entry in such a benchmark is a
\emph{structure-aware rule-based controller}: one that reads the plant topology from the
metadata and otherwise encodes physics every generated building shares. It transfers by
construction, it shows which part of the score is reachable without learning, and the gap
to the organizers' plant-agnostic baselines is, as we show, large.

\paragraph{Contributions.}
(i) An analysis of the SB26 reward and plant that yields the design facts above
(Section~\ref{sec:env}). (ii) SADR, a deterministic numpy controller (\PolicyMsOurs{}\,ms per
step) built from them (Section~\ref{sec:method}). (iii) A tuning protocol -- random search on
five-day screens, an ablation, then a decomposition of the full-length loss that motivated
one revision verified at full length on all 30 public scenarios and once on \NHeld{}
scenarios from generated buildings (Section~\ref{sec:results}). (iv) An account of where the
loss sits and where SADR still loses (Section~\ref{sec:analysis}), every number generated
from the result files.

\section{Environment and metric}
\label{sec:env}

\paragraph{Plant.}
Each building is a multi-storey floor plan partitioned into zones; air handlers (AHUs) each
serve up to ten zones through one VAV box per zone. An AHU mixes return air with outside
air (recirculation fraction $r\le 0.7$) and cools the mixture to its supply-air temperature
(SAT) setpoint, so the supply is $\min(\text{mixed}, \text{SAT})$: the AHU cannot heat. Zone
flow is $d\,F_{\max}\sqrt{P/P_{\text{design}}}$ for damper position $d\in[0.1,1]$ and static
pressure $P\in[200,2000]$\,Pa (both floored by the simulator while the unit runs). Fan power
is $\text{flow}\times P/\eta$, so the cheapest pressure for a given flow is the floor and
power grows as $P^{3/2}$ once dampers are open; the chiller's coefficient of performance is a
constant 3.5. Heat comes only from each VAV's reheat coil, fed by one gas boiler loop whose
supply temperature (310--350\,K) and differential pressure (0--20\,bar) are setpoints; the
coil's effectiveness $\min(1, C_w/C_a, (308.15-T_{\text{sa}})/(T_w - T_{\text{sa}}))$ caps the
discharge at 308.15\,K. Gas is billed as heat delivered plus a tank loss.

\paragraph{Reward and metric.}
Each 5-minute step pays $0.6\,\text{(comfort shortfall)} + 0.2\,\text{(normalized energy
cost)} + 0.2\,\text{(normalized carbon)}$. The comfort shortfall is the occupancy-weighted
productivity loss (a deviation of $\Delta$\,K from the band costs $\Delta/5$ up to 5\,K and
$1+((\Delta-5)/5)^2$ beyond, normalized by total occupancy), so it is \emph{zero whenever
nobody is in the building}; occupants are present only on weekday comfort hours (arrivals
06--12, departures 13--19, none on US holidays) and only the building-wide count is
observed. The band is 294--297\,K in comfort hours (06:00--19:00 weekdays) and 289--298\,K
otherwise. Energy is normalized by a plant maximum that scales with $K$; from the kit's cost
tables one kW of electricity costs about $2.1\times$ one kW of gas in metric units and one kW
of chiller cooling about $0.6\times$. The leaderboard metric is the mean over scenarios of
$|\text{mean per-step reward}|$, lower is better; it is dominated by the extremes: the
organizers' \texttt{pid\_water} loses \ExtremeRatioFairbanks$\times$ (Fairbanks) and
\ExtremeRatioPhoenix$\times$ (Phoenix) its mean loss in the three mild climates.

\paragraph{Baselines.}
The kit ships \texttt{bang\_bang}, \texttt{fixed\_setpoint}, \texttt{pid} (per-AHU PID on the
mean served-zone temperature, pressure scaled with effort, plant idle outside work hours) and
\texttt{pid\_water} (\texttt{pid} plus worst-zone tracking, a proportional economizer,
proportional reheat and a demand-based hot-water reset). \texttt{pid\_water} is the strongest
and is our main comparison; there is no learned baseline.

\section{Method: SADR}
\label{sec:method}

SADR is a single \texttt{TransferPolicy} class whose only state is per-zone integrals,
per-AHU dwell counters and the optimum-start statistics; all gains live in one parameter
dictionary (Table~\ref{tab:params}); Algorithm~\ref{alg:sadr} (Appendix~\ref{app:figs}) summarizes one step.

\paragraph{Topology from metadata.}
At construction the controller builds index tables from
\texttt{observation\_field\_names}, the positional AHU action lists, the per-VAV
\texttt{zone\_action\_indices} and \texttt{ahu\_zone\_assignments}, giving a zone-to-AHU
membership matrix; native action ranges come from \texttt{action\_normalization\_metadata}
and zone temperatures are de-normalized with the observation statistics. Every action slot
is written on every step and clipped to $[-1,1]$; a zone whose sensor or AHU cannot be placed
gets its damper and valve driven to the floor.

\paragraph{Modes.}
\textsc{Occupied} when the \texttt{comfort\_mode\_now} flag is set and the building is not
detected as closed; \textsc{Precondition} when unoccupied but comfort hours start within a
lead time (or \texttt{comfort\_mode\_soon}); \textsc{Setback} otherwise. The occupancy count
is affinely normalized with an unknown constant, so ``somebody is here'' is detected
relative to the running minimum of the signal; a weekday on which nobody has arrived by
\ParHolidayCheckH{}:00 is a holiday in the kit's occupancy model and is treated as
unoccupied until the count rises. The optimum-start lead is $\ell = \ell_0 +
\text{deficit}/(\hat{r}/s)$: the worst zone's distance from the comfort band over a
recovery-rate estimate $\hat{r}$ (K/h) re-estimated from each morning's recovery
(exponential average, weight \ParRateEma) with safety factor $s=\ParLeadSafety$, clipped to
$[\ParLeadMinH, \ParLeadMaxH]$\,h.

\paragraph{Zone loops on net heat.}
Each zone has heating and cooling targets, the sensor band shrunk by \ParMarginOcc\,K
(occupied) or \ParMarginPre\,K (preconditioning). Two PI loops compute requested net heat
$q_h = k_q (T_h^\star - T)^+ + I_h$ and cooling $q_c = k_q (T - T_c^\star)^+ + I_c$ in watts
($k_q = \ParKqWPerK$\,W/K, integral gain \ParKiWPerKStep\,W/K per step); the integrals unwind
past the target, so a zone inside the band free-floats, and the larger request wins. Net
heat rather than flow because the zone-air node responds within minutes: a flow-proportional
controller with a 5-minute step oscillates (our first version swung zone temperatures across
the whole eco band in one step at full pressure), whereas a net-heat gain of the order of
the air node's conductance to the slab is stable and the integral carries the steady load.
\textsc{Setback} uses \emph{bursts} instead: a running handler pays the ventilation load of
every zone it serves, so the floor (eco band widened by \ParNightLowExt{}/\ParNightHighExt\,K)
is held by requesting \ParQSbW\,W from any zone that crosses it until it is \ParSbHystK\,K
inside, and the handler is off otherwise. A version that held the floor with the PI loops
kept at least one handler running for \VthreeAhuOnUnocc{} of the unoccupied steps on the
coldest public scenario, mostly heating outside air at floor flow.

\paragraph{From heat to actuators.}
A cooling zone needs capacity rate $F = q_c/(T - T_{\text{sa}})$, $T_{\text{sa}}$ being the
AHU's estimated supply temperature; a heating zone is bounded by the discharge limit, $F =
q_h/(308.15 - \delta - T)$. \emph{SAT:} a cooling handler commands the coldest value,
\ParSatCoolHi\,K. The coil load is $F\,(T_{\text{mixed}} - T_{\text{sa}}) = q_c + F\,
(T_{\text{mixed}} - T)$ and the COP is constant, so for a given $q_c$ both the ventilation term
and the fan power fall with flow and the coldest supply is cheapest; v4 instead reset SAT
upward to the warmest value that still served the worst zone at floor pressure, in the
trim-and-respond spirit of Guideline 36~\citep{ashrae2021g36,taylor2018vav} -- right in a real
building (part-load reheat, humidity), wrong here (Section~\ref{sec:analysis}). A handler
that also heats is floored at \ParSatMixFloor\,K so its cooling zones are not overcooled into
reheat; a heating-only handler commands 305\,K so the coil idles. \emph{Static pressure} is
trim-and-respond in closed form, $P = P_{\text{design}}\max_z f_z^2$ clipped to
$[200, 2000]$\,Pa with $f_z$ each zone's requested fraction of design flow, so the most
demanding damper is fully open and pressure leaves the floor only when it must; dampers are
$f_z/\sqrt{P/P_{\text{design}}}$. \emph{Reheat:} the discharge temperature that delivers $q_h$
at the delivered flow gives the water capacity rate $W_z = F_z\,(T_{\text{disch}} -
T_{\text{sa}})/(T_w - T_{\text{sa}})$; the loop differential pressure is reset so the largest
$W_z$ sits at \ParValveTarget{} of valve travel and each valve is $W_z/(k_v \sqrt{\Delta p})$
with the kit's $k_v$. An AHU runs iff a served zone requests more than \texttt{q\_on}
(dwell \ParMinOnSteps{} steps); the pump runs iff a valve is open. Recirculation is at the 0.7
cap except when a cooling-only handler sees outside air colder than return air by more than
the exhaust-fan break-even ($0.00322$\,K/Pa$\times P$), in which case the outside-air
fraction is solved so the mixed air lands on the SAT target. The hot-water setpoint is 350\,K
whenever heating was requested in the last \ParHeatSeasonSteps{} steps and 310\,K otherwise:
gas is billed as heat delivered and the discharge limit makes hotter water strictly more
useful, at an extra tank dissipation of \TankLossExtraKw\,kW by the kit's boiler model.


\paragraph{Transfer safeguards.}
Nothing depends on the number of zones or handlers, the floor plan or the strip; the only
building-specific constants are read from metadata or observations (design flow per VAV
from \texttt{supply\_air\_flowrate\_setpoint}). $k_v$ and the discharge limit are kit
constants; if a held-out building changed them the PI loops would still regulate with a
mis-scaled proportional path. Non-finite inputs fall back to schedule values; no files, no
network, no randomness.

\paragraph{Tuning protocol.}
Parameters were first tuned by seeded random search on a screening set of six scenarios
from four public buildings (1, 2, 5, 8; one strip per station plus a second Fairbanks strip)
at 1,440 steps (five days): the design defaults plus \NTuneRandom{} random perturbations of
two to five parameters (\NTuneConfigs{} distinct configurations). The defaults scored
\TuneDefaultScreen, the best \TuneBestScreen{} and \texttt{pid\_water} \TunePidWaterScreen{}
on the same runs. The best configuration (v4) differs from the defaults in four values, and
a one-at-a-time ablation (Appendix~\ref{app:tuning}) shows that doubling the integral gain to
\ParKiWPerKStep\,W/K per step is the whole gain (\AblRelDefAddKi{} on its own; the other three
are within 1\%). Buildings 4 and 7 were never screened. v4 was then run on all 30 public
scenarios at 8,640 steps, its loss decomposed (Section~\ref{sec:analysis}), and \NCandV{}
physics-motivated candidates (cold SAT; night setback widened to 8--12\,K; combinations) were
screened the same way and all run at full length on the 30 public scenarios
(Appendix~\ref{app:tuning}). The winner on the public set is the final policy (v5); the
\NHeld{} generated-building scenarios were run once, after that choice.

\begin{table}[t]
\caption{Final SADR parameters (v5); values that changed from v4 are marked. The full dictionary is in the code.}
\label{tab:params}
\centering\small
\begin{tabular}{llll}
\toprule
occupied margin & \ParMarginOcc\,K & precondition margin & \ParMarginPre\,K \\
setback widening (low/high) & \ParNightLowExt{}/\ParNightHighExt\,K (v4: \VFourNightLowExt/\VFourNightHighExt) & $k_q$, $k_i$ & \ParKqWPerK, \ParKiWPerKStep\,W/K \\
cooling SAT / mixed floor & \ParSatCoolHi{}/\ParSatMixFloor\,K (v4 ceiling: \VFourSatCoolHi) & discharge margin & \ParDischMarginK\,K \\
lead safety / rate$_0$ & \ParLeadSafety{} / \ParRateZeroKph\,K/h & lead range & \ParLeadMinH--\ParLeadMaxH\,h \\
heating-season memory & \ParHeatSeasonSteps{} steps & AHU dwell & \ParMinOnSteps{} steps \\
\bottomrule
\end{tabular}
\end{table}

\section{Results}
\label{sec:results}

\paragraph{Setup.}
All runs use the kit's environment loader and \texttt{SimulationTracker}, 8,640 steps per
scenario, and the kit's \texttt{score\_leaderboard.py} formula with paired bootstrap CIs over
scenarios added. Every policy in Table~\ref{tab:main} -- SADR v5, SADR v4 and the four
organizer baselines -- was run on one machine (Apple M5 Pro, \SecPerStepOurs\,s per simulator
step with 14 concurrent runs); local \texttt{pid\_water} returns are \GapMeanPidWater{}
relative to the shipped \texttt{baseline\_returns.json} on average (range \GapMinPidWater{} to
\GapMaxPidWater), so all paired comparisons use local runs. As a cross-machine check, v4 and
\texttt{pid\_water} were also run on a Windows laptop (\SimSecPerStep\,s per step): v4 scores
\LocalVFourMetric{} there against \MetricOursVFour{} on the Mac, a per-scenario relative
difference of \MachineGapMean{} (max \MachineGapMax; \MachineGapPidWater{} for
\texttt{pid\_water}); the simulator's seeds are fixed, so the numbers reproduce across machines.

\begin{table}[t]
\caption{Leaderboard metric (mean over the \NScen{} public scenarios of $|$mean per-step
reward$|$; lower is better) and per-station means (diagnostic only; the leaderboard never
groups). Paired bootstrap 95\% CIs (\NBoot{} resamples over scenarios) are for SADR v5 minus
the row (negative favours v5). All rows were run on one machine; the last row is the
organizers' shipped returns. Per-scenario numbers: Appendix~\ref{app:per_scenario}.}
\label{tab:main}
\centering\scriptsize
\setlength{\tabcolsep}{3pt}
\begin{tabular}{lcccccc c}
\toprule
policy & all & Fairbanks & Denver & Atlanta & Phoenix & Singapore & $\Delta$ vs v5 [95\% CI] \\
\midrule
SADR v5 (final) & \textbf{\MetricOurs} & \MetricOursFairbanks & \MetricOursDenver & \MetricOursAtlanta & \MetricOursPhoenix & \MetricOursSingapore & -- \\
SADR v4 & \MetricOursVFour & \MetricOursVFourFairbanks & \MetricOursVFourDenver & \MetricOursVFourAtlanta & \MetricOursVFourPhoenix & \MetricOursVFourSingapore & \DiffOursVFour{} [\DiffLoOursVFour, \DiffHiOursVFour] \\
pid\_water & \MetricPidWater & \MetricPidWaterFairbanks & \MetricPidWaterDenver & \MetricPidWaterAtlanta & \MetricPidWaterPhoenix & \MetricPidWaterSingapore & \DiffPidWater{} [\DiffLoPidWater, \DiffHiPidWater] \\
pid & \MetricPid & \MetricPidFairbanks & \MetricPidDenver & \MetricPidAtlanta & \MetricPidPhoenix & \MetricPidSingapore & \DiffPid{} [\DiffLoPid, \DiffHiPid] \\
fixed\_setpoint & \MetricFixedSp & \MetricFixedSpFairbanks & \MetricFixedSpDenver & \MetricFixedSpAtlanta & \MetricFixedSpPhoenix & \MetricFixedSpSingapore & \DiffFixedSp{} [\DiffLoFixedSp, \DiffHiFixedSp] \\
bang\_bang & \MetricBangBang & \MetricBangBangFairbanks & \MetricBangBangDenver & \MetricBangBangAtlanta & \MetricBangBangPhoenix & \MetricBangBangSingapore & \DiffBangBang{} [\DiffLoBangBang, \DiffHiBangBang] \\
\midrule
pid\_water (organizer values) & \MetricPidWaterOff & \MetricPidWaterOffFairbanks & \MetricPidWaterOffDenver & \MetricPidWaterOffAtlanta & \MetricPidWaterOffPhoenix & \MetricPidWaterOffSingapore & \DiffPidWaterOff{} [\DiffLoPidWaterOff, \DiffHiPidWaterOff] \\
\bottomrule
\end{tabular}
\end{table}

\paragraph{Main result.}
SADR v5 reaches \MetricOurs{} against \MetricPidWater{} for \texttt{pid\_water}
(Table~\ref{tab:main}, Figure~\ref{fig:scores}), a relative reduction of \RelRedPidWater{}
(95\% CI \RelRedLoPidWater{} to \RelRedHiPidWater), winning on \WinsPidWater{} of \NScen{}
scenarios (Wilcoxon signed-rank $p\PWilcoxonPidWater$); against the shipped organizer values
the ratio is \RatioPidWaterOff. The absolute reduction is largest in Fairbanks
(\DiffPidWaterFairbanks) and Phoenix (\DiffPidWaterPhoenix). The revision from v4 to v5 cut
the metric by \RelRedOursVFour{} (CI \RelRedLoOursVFour{} to \RelRedHiOursVFour; wins
\WinsOursVFour/\NScen), concentrated in the cooling climates (Phoenix \DiffOursVFourPhoenix,
Singapore \DiffOursVFourSingapore); Fairbanks moves by \DiffOursVFourFairbanks{} and the mild
climates by at most \DiffOursVFourAtlanta. For scale, the
public Kaggle leaderboard (read \LbDate; \LbNTeams{} teams; a third of the held-out test
data; same formula on unseen buildings and strips, so not directly comparable) stood at
\LbRankOne{} for rank 1, \LbRankEight{} for rank 8 (the last prize rank), \LbRankThirty{} for
rank 30 and \LbRankFortyNine{} for rank 49, against \MetricOurs{} here
(\OursVsLbOne$\times$ the leader).

\begin{figure}[t]
\centering
\figfile{\linewidth}{figures/fig_scores_by_scenario.pdf}
\caption{Per-scenario leaderboard score (log scale, lower is better), SADR v5 vs
\texttt{pid\_water} on the same machine, grouped by station and labelled by building (b =
public, g = generated held-out). Left: the 30 public scenarios; right: the \NHeld{}
held-out scenarios.}
\label{fig:scores}
\end{figure}

\paragraph{Where the loss goes.}
Figure~\ref{fig:decomp} splits each policy's per-step loss into comfort, energy and carbon
per station, and Figure~\ref{fig:hours} by hour of day. Over the \NAnalysis{} public
scenarios, comfort regret is \RegretShareOurs{} of SADR v5's loss (energy \EnergyShareOurs,
carbon \CarbonShareOurs) against \RegretSharePidWater{} of \texttt{pid\_water}'s (energy
\EnergySharePidWater, carbon \CarbonSharePidWater). A one-day trace of both controllers on
the coldest scenario is in Appendix~\ref{app:figs}.

\begin{figure}[t]
\centering
\figfile{0.58\linewidth}{figures/fig_station_decomp.pdf}
\caption{Per-station decomposition of the per-step loss (metric units) into comfort,
energy and carbon, SADR v5 vs \texttt{pid\_water}.}
\label{fig:decomp}
\end{figure}

\begin{figure}[t]
\centering
\figfile{\linewidth}{figures/fig_loss_by_hour.pdf}
\caption{Mean per-step loss by hour of day over the \NAnalysis{} public scenarios, split
into comfort regret and energy\,+\,carbon; comfort hours (06--19) shaded. Left: SADR v5;
right: \texttt{pid\_water}.}
\label{fig:hours}
\end{figure}

\paragraph{Held-out buildings.}
On the \NHeld{} scenarios of the two generated buildings (\ZonesBgena{} zones/\AhusBgena{}
AHUs and \ZonesBgenb{} zones/\AhusBgenb{} AHUs; strips absent from the public set), run once
after the final policy was chosen, SADR v5 scores \HeldMetricOurs{} against
\HeldMetricPidWater{} for \texttt{pid\_water}, a ratio of \HeldRatio{} (difference \HeldDiff,
95\% CI [\HeldDiffLo, \HeldDiffHi]; wins \HeldWins/\NHeld): the ratio carries over, less
favourably mainly on the larger generated building. The v4-to-v5 gain shrinks to
\HeldVFourDiff{} (v4 \HeldMetricOursVFour; ratio \HeldVFourRatio, CI [\HeldVFourDiffLo,
\HeldVFourDiffHi], wins \HeldVFourWins/\NHeld) and is not significant on ten scenarios. The
public buildings never screened (4 and 7) behave like the others (v5 \MetricOursBfour{} and
\MetricOursBseven{} vs \MetricPidWaterBfour{} and \MetricPidWaterBseven).

\paragraph{Tuning and candidates.}
The spread of the \NTuneConfigs{} random-search configurations on the screen
(\TuneBestScreen{} to \TuneWorstScreen, Appendix~\ref{app:figs}) is small relative to the gap
to \texttt{pid\_water} (\TunePidWaterScreen{} on the same runs): the gains come from the
structure, not the tuning. The \NCandV{} analysis-driven candidates score \CandBestScreen{} at
best on the screen (cold SAT alone \CandSatColdScreen; the deeper setback alone
\CandNightEightScreen, which a five-day screen cannot see); at full length the winner,
\CandWinner, reaches \CandWinnerMetric{} against \CandWinnerVFourMetric{} for v4 on the same
machine (ratio \CandWinnerRatio, CI [\CandWinnerCiLo, \CandWinnerCiHi], wins
\CandWinnerWins/\NScen); all candidates are in Appendix~\ref{app:tuning}.

\section{Analysis}
\label{sec:analysis}

\paragraph{Why it wins.}
The decomposition is unambiguous about the baseline: for \texttt{pid\_water} comfort regret
is \RegretSharePidWater{} of the loss and \OccOutPidWater{} of occupied zone-steps lie outside
the comfort band (mean worst-zone deviation \OccWorstDevPidWater\,K), because its handler
loop regulates one temperature per handler while individual zones drift, and its start is a
step at 06:00 (the spike in Figure~\ref{fig:hours}). SADR keeps \OccOutOurs{} of occupied
zone-steps out of band (comfort \RegretShareOurs{} of its loss) with zone-level integral
control and a learned lead. In the extreme climates it also uses far less energy: pressure
at the floor with wide dampers (ten times less fan power than the 2000\,Pa
\texttt{pid\_water} commands on a 2\,K excursion, for the same flow), hot water at 350\,K with
the minimum flow the discharge cap allows (the ventilation penalty per joule of heat,
$(T - T_{\text{mixed}})/(T_{\text{disch}} - T)$, is minimized), recirculation at the cap, and no
conditioning while nobody is charged for comfort.

\paragraph{Where the loss sits now, and what the revision bought.}
For SADR v5, \UnoccShareOurs{} of the loss is incurred on unoccupied steps
(\UnoccShareOursFairbanks{} in Fairbanks), and the largest hourly spend is hour
\PeakLossHourOurs{} (\PeakLossShareOurs{} of the day), the warm-up and pre-cool before 06:00. Fairbanks is gas: heating with a mandatory 30\%
outside-air fraction at $-20^\circ$C costs about $1.9$\,J of gas per joule of net heat at
the discharge cap, and no plant in this benchmark heats any other way; what remains
adjustable is the envelope loss, hence the deeper setback. Phoenix and Singapore were
electricity for cooling, and the v4 SAT reset was spending flow: with coil load $q_c +
F\,(T_{\text{mixed}} - T)$ and a constant COP, the coldest supply at the lowest flow is
strictly cheaper. The candidates test exactly these two levers and the full-length
differences (Table~\ref{tab:main}, Appendix~\ref{app:tuning}) match the prediction: the
cooling climates improve with cold SAT, Fairbanks improves slightly with the deeper setback
(cold SAT costs it a little on mixed handlers, so the two nearly cancel there), the mild
climates barely move. The first workday, before any recovery rate has been measured, is not a weak
spot either: its comfort regret exceeds the mean of the later workdays on
\NFirstWorkdayWorse{} of \NFirstWorkdayTotal{} scenarios (median ratio
\FirstWorkdayRatioMedianOurs), and the comfort term is \RegretShareOurs{} of the loss in any
case.

\paragraph{Where it loses.}
SADR v5 is worse than \texttt{pid\_water} on \NScenWorse{} of \NScen{} scenarios; its
largest residual loss is on \MaxOursScen{} (\MaxOursScore). The failure modes we observe are
(i) in the coldest strips the morning recovery still leaves regret in the first two comfort
hours (Figure~\ref{fig:hours}): the lead is tuned to be cheap, not perfect; (ii) in the
hottest strips (outside air up to \PhoenixPeakOatC\,$^\circ$C) perimeter zones exceed what
floor pressure can serve at the afternoon peak and the pressure escalation costs fan energy;
(iii) on mixed handlers in shoulder weather the \ParSatMixFloor\,K floor heats some zones with
gas while cooling their neighbours. None is a transfer failure: the same modes appear on
screened, unscreened and generated buildings alike.

\paragraph{What this says about foundation-model control.}
The metric rewards never letting an occupied zone drift and never moving more air, at higher
pressure or warmer than necessary; both are reachable from the metadata and the physics
alone, and one analysis-driven revision was worth more than the whole parameter search. A
learned policy should be compared with a controller of this kind and its advantage looked
for where rules are weak: the first morning, the afternoon peak, shared handlers.

\section{Limitations}
SADR encodes constants of this simulator (the discharge limit, the valve coefficient, the
constant COP, the exhaust-fan break-even); in a real building these would be commissioning
parameters, and the cold-SAT rule would have to be weighed against humidity and reheat. The
candidate selection used the 30 public scenarios at full length, so the public-set number is
optimistic; the \NHeld{} held-out scenarios, run once on two generated buildings, are the
honest one, and the Kaggle test buildings may differ in ways the generator's public ranges do
not show. The controller models neither the tank's nor the slab's thermal inertia; a model
predictive controller or a learned residual on top of SADR is the natural next step.

\section{Reproducibility and disclosure}
Code: \url{https://github.com/neelsirivara/neurips2026-smart-buildings} (kit and time series
excluded). \texttt{submission/transfer\_policy.py} is the policy (drop-in for the kit);
\texttt{experiments/src/} holds the resumable evaluators (\texttt{run\_scenario.py},
\texttt{run\_eval\_parallel.py}), the search (\texttt{tune.py}, \texttt{ablation.py},
\texttt{compare\_candidates.py}), the scoring (\texttt{score\_all.py},
\texttt{analyze\_final.py}) and the generators of every number and figure here
(\texttt{make\_macros.py}, \texttt{make\_figures.py}); the run scripts are
\texttt{mac\_jobs.sh}, \texttt{mac\_final.sh} and \texttt{run\_final\_eval.sh}. The simulator's
occupancy (17321) and convection (5) seeds are fixed in the kit; the bootstrap uses seed 0.
Python 3.12.15 with the kit's pins on both machines; about \TotalCoreHours{} core-hours on the
laptop plus the Mac runs. Compute was a CPU-only Windows laptop (8 cores, 34\,GB RAM) and an
Apple M5 Pro (18 cores, 48\,GB RAM), no GPU. An LLM coding assistant (Claude) was used to write
code and to draft and edit this text; the design decisions and the verification of every
number against the result files are recorded in \texttt{paper/AUDIT.md}. No external data or
models beyond the organizers' starter kit were used.

\bibliographystyle{plainnat}
\bibliography{refs}

\appendix
\section{Per-scenario results}
\label{app:per_scenario}
Tables~\ref{tab:per_scenario} and~\ref{tab:per_scenario_held} list every scenario's total
return and leaderboard score; all values are generated from the result files by
\texttt{make\_macros.py}.
% AUTO-GENERATED by experiments/src/make_macros.py -- do not edit by hand
\begin{table}[!htbp]
\caption{Per-scenario results on the 30 public scenarios (8,640 steps). Score is $|$total$/8640|$, lower is better; bold marks scenarios where SADR beats the locally re-run \texttt{pid\_water}; organizer = the \texttt{pid\_water} score from the kit's \texttt{baseline\_returns.json}.}
\label{tab:per_scenario}
\centering\scriptsize
\setlength{\tabcolsep}{2.5pt}
\begin{tabular}{llrrrrrrr}
\toprule
scenario & zones/AHUs & SADR v6 return & v6 score & v5 score & v4 score & pid\_water return & pid\_water score & organizer \\
\midrule
building\_train\_1/atlanta\_2 & 39/4 & -34.3 & \textbf{0.0040} & 0.0041 & 0.0037 & -215.2 & 0.0249 & 0.0241 \\
building\_train\_2/atlanta\_3 & 85/9 & -29.2 & \textbf{0.0034} & 0.0036 & 0.0036 & -172.6 & 0.0200 & 0.0200 \\
building\_train\_4/atlanta\_5 & 54/6 & -58.3 & \textbf{0.0067} & 0.0069 & 0.0085 & -250.0 & 0.0289 & 0.0285 \\
building\_train\_5/atlanta\_6 & 79/8 & -38.0 & \textbf{0.0044} & 0.0044 & 0.0055 & -265.8 & 0.0308 & 0.0273 \\
building\_train\_7/atlanta\_8 & 36/4 & -66.6 & \textbf{0.0077} & 0.0080 & 0.0080 & -417.3 & 0.0483 & 0.0483 \\
building\_train\_8/atlanta\_1 & 30/3 & -48.3 & \textbf{0.0056} & 0.0059 & 0.0055 & -267.4 & 0.0309 & 0.0325 \\
building\_train\_1/denver\_7 & 39/4 & -24.6 & \textbf{0.0028} & 0.0029 & 0.0025 & -267.9 & 0.0310 & 0.0299 \\
building\_train\_2/denver\_8 & 85/9 & -43.0 & \textbf{0.0050} & 0.0055 & 0.0055 & -351.1 & 0.0406 & 0.0398 \\
building\_train\_4/denver\_2 & 54/6 & -53.8 & \textbf{0.0062} & 0.0064 & 0.0063 & -471.4 & 0.0546 & 0.0545 \\
building\_train\_5/denver\_3 & 79/8 & -40.5 & \textbf{0.0047} & 0.0047 & 0.0044 & -381.0 & 0.0441 & 0.0437 \\
building\_train\_7/denver\_5 & 36/4 & -50.5 & \textbf{0.0058} & 0.0061 & 0.0068 & -318.6 & 0.0369 & 0.0361 \\
building\_train\_8/denver\_6 & 30/3 & -33.8 & \textbf{0.0039} & 0.0041 & 0.0046 & -176.9 & 0.0205 & 0.0202 \\
building\_train\_1/fairbanks\_1 & 39/4 & -420.3 & \textbf{0.0487} & 0.0506 & 0.0514 & -1849.7 & 0.2141 & 0.2117 \\
building\_train\_2/fairbanks\_2 & 85/9 & -111.4 & \textbf{0.0129} & 0.0141 & 0.0150 & -684.6 & 0.0792 & 0.0798 \\
building\_train\_4/fairbanks\_4 & 54/6 & -21.0 & \textbf{0.0024} & 0.0026 & 0.0019 & -283.4 & 0.0328 & 0.0316 \\
building\_train\_5/fairbanks\_5 & 79/8 & -20.0 & \textbf{0.0023} & 0.0023 & 0.0019 & -145.9 & 0.0169 & 0.0168 \\
building\_train\_7/fairbanks\_7 & 36/4 & -75.6 & \textbf{0.0087} & 0.0091 & 0.0091 & -467.1 & 0.0541 & 0.0537 \\
building\_train\_8/fairbanks\_8 & 30/3 & -218.8 & \textbf{0.0253} & 0.0264 & 0.0278 & -668.7 & 0.0774 & 0.0772 \\
building\_train\_1/phoenix\_5 & 39/4 & -129.2 & \textbf{0.0150} & 0.0153 & 0.0198 & -1183.8 & 0.1370 & 0.1335 \\
building\_train\_2/phoenix\_6 & 85/9 & -64.2 & \textbf{0.0074} & 0.0075 & 0.0105 & -647.9 & 0.0750 & 0.0708 \\
building\_train\_4/phoenix\_8 & 54/6 & -20.8 & \textbf{0.0024} & 0.0025 & 0.0020 & -261.5 & 0.0303 & 0.0288 \\
building\_train\_5/phoenix\_1 & 79/8 & -27.9 & \textbf{0.0032} & 0.0033 & 0.0027 & -287.4 & 0.0333 & 0.0333 \\
building\_train\_7/phoenix\_3 & 36/4 & -55.2 & \textbf{0.0064} & 0.0065 & 0.0078 & -387.3 & 0.0448 & 0.0444 \\
building\_train\_8/phoenix\_4 & 30/3 & -87.3 & \textbf{0.0101} & 0.0102 & 0.0145 & -1384.2 & 0.1602 & 0.1580 \\
building\_train\_1/singapore\_8 & 39/4 & -57.5 & \textbf{0.0067} & 0.0068 & 0.0085 & -244.9 & 0.0284 & 0.0282 \\
building\_train\_2/singapore\_1 & 85/9 & -41.5 & \textbf{0.0048} & 0.0049 & 0.0063 & -223.8 & 0.0259 & 0.0256 \\
building\_train\_4/singapore\_3 & 54/6 & -60.9 & \textbf{0.0071} & 0.0075 & 0.0095 & -296.3 & 0.0343 & 0.0325 \\
building\_train\_5/singapore\_4 & 79/8 & -48.5 & \textbf{0.0056} & 0.0056 & 0.0076 & -949.6 & 0.1099 & 0.1085 \\
building\_train\_7/singapore\_6 & 36/4 & -50.7 & \textbf{0.0059} & 0.0058 & 0.0077 & -260.9 & 0.0302 & 0.0311 \\
building\_train\_8/singapore\_7 & 30/3 & -48.8 & \textbf{0.0057} & 0.0057 & 0.0074 & -468.2 & 0.0542 & 0.0541 \\
\bottomrule
\end{tabular}
\end{table}

\begin{table}[!htbp]
\caption{Per-scenario results on the generated held-out buildings 0501 and 0502 (8,640 steps), SADR vs locally run \texttt{pid\_water}.}
\label{tab:per_scenario_held}
\centering\scriptsize
\setlength{\tabcolsep}{2.5pt}
\begin{tabular}{llrrrrrr}
\toprule
scenario & zones/AHUs & SADR v6 return & v6 score & v5 score & v4 score & pid\_water return & pid\_water score \\
\midrule
building\_0501/atlanta\_7 & 29/5 & -32.8 & \textbf{0.0038} & 0.0043 & 0.0042 & -241.6 & 0.0280 \\
building\_0502/atlanta\_4 & 28/4 & -77.2 & \textbf{0.0089} & 0.0099 & 0.0107 & -305.1 & 0.0353 \\
building\_0501/denver\_4 & 29/5 & -32.8 & \textbf{0.0038} & 0.0044 & 0.0044 & -226.4 & 0.0262 \\
building\_0502/denver\_1 & 28/4 & -228.8 & \textbf{0.0265} & 0.0325 & 0.0325 & -616.1 & 0.0713 \\
building\_0501/fairbanks\_6 & 29/5 & -24.2 & \textbf{0.0028} & 0.0030 & 0.0027 & -291.8 & 0.0338 \\
building\_0502/fairbanks\_3 & 28/4 & -202.6 & \textbf{0.0234} & 0.0285 & 0.0280 & -598.6 & 0.0693 \\
building\_0501/phoenix\_7 & 29/5 & -31.5 & \textbf{0.0036} & 0.0040 & 0.0044 & -232.0 & 0.0269 \\
building\_0502/phoenix\_2 & 28/4 & -44.3 & \textbf{0.0051} & 0.0071 & 0.0069 & -240.0 & 0.0278 \\
building\_0501/singapore\_5 & 29/5 & -50.0 & \textbf{0.0058} & 0.0064 & 0.0075 & -345.7 & 0.0400 \\
building\_0502/singapore\_2 & 28/4 & -91.2 & \textbf{0.0106} & 0.0107 & 0.0123 & -348.7 & 0.0404 \\
\bottomrule
\end{tabular}
\end{table}

\begin{table}[!htbp]
\caption{Per-scenario results on the six generated held-out buildings of round 3 (8,640 steps), SADR vs locally run \texttt{pid\_water}.}
\label{tab:per_scenario_held_r3}
\centering\scriptsize
\setlength{\tabcolsep}{2.5pt}
\begin{tabular}{llrrrrrr}
\toprule
scenario & zones/AHUs & SADR v6 return & v6 score & v5 score & v4 score & pid\_water return & pid\_water score \\
\midrule
building\_0500/atlanta\_4 & 120/12 & -104.2 & \textbf{0.0121} & 0.0130 & -- & -344.0 & 0.0398 \\
building\_0505/atlanta\_1 & 63/7 & -51.2 & \textbf{0.0059} & 0.0069 & -- & -525.2 & 0.0608 \\
building\_0507/atlanta\_2 & 68/7 & -83.2 & \textbf{0.0096} & 0.0113 & -- & -409.9 & 0.0474 \\
building\_0508/atlanta\_7 & 89/9 & -66.2 & \textbf{0.0077} & 0.0090 & -- & -334.3 & 0.0387 \\
building\_0509/atlanta\_3 & 32/5 & -62.3 & \textbf{0.0072} & 0.0084 & -- & -302.2 & 0.0350 \\
building\_0511/atlanta\_5 & 71/8 & -64.2 & \textbf{0.0074} & 0.0075 & -- & -279.4 & 0.0323 \\
building\_0500/denver\_1 & 120/12 & -197.5 & \textbf{0.0229} & 0.0274 & -- & -828.1 & 0.0958 \\
building\_0505/denver\_2 & 63/7 & -87.3 & \textbf{0.0101} & 0.0114 & -- & -755.0 & 0.0874 \\
building\_0507/denver\_3 & 68/7 & -76.4 & \textbf{0.0088} & 0.0097 & -- & -392.7 & 0.0455 \\
building\_0508/denver\_4 & 89/9 & -69.3 & \textbf{0.0080} & 0.0097 & -- & -310.3 & 0.0359 \\
building\_0509/denver\_5 & 32/5 & -82.5 & \textbf{0.0095} & 0.0107 & -- & -330.7 & 0.0383 \\
building\_0511/denver\_6 & 71/8 & -40.7 & \textbf{0.0047} & 0.0049 & -- & -219.4 & 0.0254 \\
building\_0500/fairbanks\_3 & 120/12 & -175.8 & \textbf{0.0203} & 0.0242 & -- & -843.3 & 0.0976 \\
building\_0505/fairbanks\_1 & 63/7 & -435.6 & \textbf{0.0504} & 0.0568 & -- & -1763.8 & 0.2041 \\
building\_0507/fairbanks\_2 & 68/7 & -213.1 & \textbf{0.0247} & 0.0304 & -- & -1064.4 & 0.1232 \\
building\_0508/fairbanks\_6 & 89/9 & -44.7 & \textbf{0.0052} & 0.0062 & -- & -313.7 & 0.0363 \\
building\_0509/fairbanks\_4 & 32/5 & -27.9 & \textbf{0.0032} & 0.0044 & -- & -256.1 & 0.0296 \\
building\_0511/fairbanks\_5 & 71/8 & -29.6 & \textbf{0.0034} & 0.0035 & -- & -203.2 & 0.0235 \\
building\_0500/phoenix\_2 & 120/12 & -66.3 & \textbf{0.0077} & 0.0089 & -- & -311.2 & 0.0360 \\
building\_0505/phoenix\_3 & 63/7 & -85.7 & \textbf{0.0099} & 0.0105 & -- & -323.4 & 0.0374 \\
building\_0507/phoenix\_4 & 68/7 & -272.7 & \textbf{0.0316} & 0.0326 & -- & -625.3 & 0.0724 \\
building\_0508/phoenix\_7 & 89/9 & -84.3 & \textbf{0.0098} & 0.0106 & -- & -318.8 & 0.0369 \\
building\_0509/phoenix\_5 & 32/5 & -161.2 & \textbf{0.0187} & 0.0246 & -- & -477.0 & 0.0552 \\
building\_0511/phoenix\_6 & 71/8 & -100.0 & \textbf{0.0116} & 0.0119 & -- & -452.7 & 0.0524 \\
building\_0500/singapore\_2 & 120/12 & -116.9 & \textbf{0.0135} & 0.0151 & -- & -383.9 & 0.0444 \\
building\_0505/singapore\_4 & 63/7 & -89.0 & \textbf{0.0103} & 0.0110 & -- & -312.4 & 0.0362 \\
building\_0507/singapore\_6 & 68/7 & -191.0 & \textbf{0.0221} & 0.0229 & -- & -500.3 & 0.0579 \\
building\_0508/singapore\_5 & 89/9 & -123.0 & \textbf{0.0142} & 0.0152 & -- & -386.5 & 0.0447 \\
building\_0509/singapore\_7 & 32/5 & -95.2 & \textbf{0.0110} & 0.0122 & -- & -348.7 & 0.0404 \\
building\_0511/singapore\_8 & 71/8 & -56.6 & \textbf{0.0066} & 0.0067 & -- & -253.9 & 0.0294 \\
\bottomrule
\end{tabular}
\end{table}



\section{Ablation and candidates}
\label{app:tuning}
Table~\ref{tab:ablation} reverts each tuned v4 value on its own (and adds each on its own to
the defaults) on the screening set; Table~\ref{tab:candidates} lists every v5 candidate at
full length on the 30 public scenarios against v4 on the same machine.
% AUTO-GENERATED by experiments/src/make_macros.py -- do not edit by hand
\begin{table}[h]
\caption{Ablation of the four tuned v4 values on the screening set (six scenarios, 1,440 steps; leaderboard metric, lower is better). Rev: the best configuration with one value reverted to its default; Add: the defaults with one tuned value added. Best 0.00849, defaults 0.00947.}
\label{tab:ablation}
\centering\small
\begin{tabular}{llrrr}
\toprule
parameter & change & score & vs best & vs defaults \\
\midrule
\texttt{ki\_w\_per\_k\_step} & Rev (60) & 0.00905 & +6.7\% & -4.4\% \\
\texttt{ki\_w\_per\_k\_step} & Add (120) & 0.00855 & +0.7\% & -9.8\% \\
\texttt{margin\_pre} & Rev (0.5) & 0.00846 & -0.3\% & -10.7\% \\
\texttt{margin\_pre} & Add (0.8) & 0.00910 & +7.2\% & -4.0\% \\
\texttt{sat\_mix\_floor} & Rev (290) & 0.00849 & +0.0\% & -10.4\% \\
\texttt{sat\_mix\_floor} & Add (292) & 0.00947 & +11.6\% & +0.0\% \\
\texttt{sb\_hyst\_k} & Rev (1.5) & 0.00850 & +0.2\% & -10.2\% \\
\texttt{sb\_hyst\_k} & Add (1) & 0.00944 & +11.2\% & -0.3\% \\
\bottomrule
\end{tabular}
\end{table}

\begin{table}[h]
\caption{v5 candidates at full length (8,640 steps) on the 30 public scenarios, run on the same machine as v4 (metric 0.00920). Overrides are applied to the v4 values (SAT = \texttt{sat\_cool\_hi} in K; nlo/nhi = \texttt{night\_low/high\_ext} in K; occ = \texttt{margin\_occ}; lead = \texttt{lead\_max\_h}); CI = paired bootstrap 95\% CI of the per-scenario difference vs v4 (negative favours the candidate); the last five columns are station means.}
\label{tab:candidates}
\centering\scriptsize
\setlength{\tabcolsep}{3pt}
\begin{tabular}{lrrcrrrrrr}
\toprule
overrides & metric & ratio & CI & wins & Fai. & Den. & Atl. & Pho. & Sin. \\
\midrule
SAT 285, nlo 8, nhi 6 & 0.00831 & 0.903 & [-0.00139, -0.00044] & 19/30 & 0.0175 & 0.0050 & 0.0055 & 0.0075 & 0.0060 \\
SAT 285, nlo 12, nhi 4, lead 10 & 0.00832 & 0.904 & [-0.00137, -0.00042] & 19/30 & 0.0175 & 0.0050 & 0.0055 & 0.0076 & 0.0060 \\
SAT 285, nlo 8, nhi 4 & 0.00833 & 0.905 & [-0.00136, -0.00043] & 19/30 & 0.0175 & 0.0050 & 0.0055 & 0.0076 & 0.0060 \\
SAT 285, nlo 8, nhi 4, occ 0.2 & 0.00840 & 0.913 & [-0.00132, -0.00034] & 18/30 & 0.0175 & 0.0051 & 0.0056 & 0.0077 & 0.0061 \\
SAT 285, nlo 8 & 0.00843 & 0.916 & [-0.00121, -0.00038] & 19/30 & 0.0175 & 0.0050 & 0.0055 & 0.0079 & 0.0062 \\
SAT 285, nhi 4 & 0.00846 & 0.919 & [-0.00127, -0.00029] & 16/30 & 0.0180 & 0.0050 & 0.0055 & 0.0076 & 0.0060 \\
SAT 285, nlo 8, occ 0.2 & 0.00851 & 0.924 & [-0.00115, -0.00028] & 18/30 & 0.0175 & 0.0052 & 0.0056 & 0.0080 & 0.0062 \\
SAT 285 & 0.00856 & 0.930 & [-0.00113, -0.00023] & 15/30 & 0.0181 & 0.0051 & 0.0055 & 0.0079 & 0.0062 \\
SAT 288 & 0.00866 & 0.941 & [-0.00088, -0.00023] & 15/30 & 0.0180 & 0.0050 & 0.0055 & 0.0083 & 0.0065 \\
nlo 12 & 0.00907 & 0.986 & [-0.00031, -0.00002] & 7/30 & 0.0173 & 0.0050 & 0.0058 & 0.0096 & 0.0078 \\
nlo 8 & 0.00908 & 0.987 & [-0.00024, -0.00002] & 7/30 & 0.0173 & 0.0049 & 0.0058 & 0.0096 & 0.0078 \\
nlo 8, lead 10 & 0.00908 & 0.987 & [-0.00024, -0.00002] & 7/30 & 0.0173 & 0.0049 & 0.0058 & 0.0096 & 0.0078 \\
none (v5 code, v4 values) & 0.00920 & 1.000 & [+0.00000, +0.00000] & 0/30 & 0.0178 & 0.0050 & 0.0058 & 0.0096 & 0.0078 \\
\bottomrule
\end{tabular}
\end{table}



\section{Algorithm and additional figures}
\label{app:figs}
\begin{algorithm}[h]
\caption{SADR, one step (vectorized over zones $z$ and handlers $k$)}
\label{alg:sadr}
\begin{algorithmic}[1]
\State mode $\gets$ \textsc{Occupied} / \textsc{Precondition} / \textsc{Setback} from flags, time, occupancy count, lead $\ell$
\State $(T_h^\star, T_c^\star)_z \gets$ band for the mode; update $I_h, I_c$; $q_h, q_c \gets$ PI requests (W), or setback bursts
\State run$_k \gets$ hysteresis on $\max_{z\in k}(q_h, q_c)$
\State SAT$_k \gets$ \ParSatCoolHi\,K if cooling (\ParSatMixFloor\,K if also heating), 305\,K otherwise; recirc$_k \gets$ economizer
\State $F_z \gets q_c/(T_z - T_{\text{sa},k})$ or $q_h/(308.15-\delta-T_z)$; $P_k \gets P_{\text{design}}\max_z f_z^2$; damper$_z \gets f_z/\sqrt{P_k/P_{\text{design}}}$
\State $W_z \gets F_z (T_{\text{disch},z} - T_{\text{sa}})/(T_w - T_{\text{sa}})$; $\Delta p \gets (\max_z W_z / (0.9 k_v))^2$; valve$_z \gets W_z/(k_v\sqrt{\Delta p})$
\State $T_w^{\text{sp}} \gets 350$\,K if heating season else 310\,K; pump $\gets$ any valve open
\State scatter all setpoints into the normalized action vectors; clip
\end{algorithmic}
\end{algorithm}
\begin{figure}[h]
\centering
\figfile{\linewidth}{figures/fig_trace_building_train_8_fairbanks_8_day8.pdf}
\caption{One weekday on \texttt{building\_train\_8/fairbanks\_8}: zone-temperature
envelope (min--max shaded, median line) against the comfort band (green wash), with the
building occupancy count shaded in grey. Left: SADR v5; right: \texttt{pid\_water}.}
\label{fig:trace}
\end{figure}
\begin{figure}[h]
\centering
\figfile{0.5\linewidth}{figures/fig_tuning.pdf}
\caption{Screening score (six scenarios, five days) of every random-search configuration,
sorted, with the design defaults highlighted and \texttt{pid\_water} on the same runs.}
\label{fig:tuning}
\end{figure}

\end{document}
